\section{Allocation as a Design Variable}
\label{sec:exponents}\label{sec:identification}
For $K=d_{\rm rot}/2$ pairs, standard RoPE uses $\omega_k=b^{-k/K}$.
Order a positive table from fast to slow, $\omega_0>\cdots>\omega_{K-1}$.
Its log frequencies admit
\begin{equation}
x_k=-\log\omega_k=a+Rz_k,\quad z_k=\frac{x_k-x_0}{x_{K-1}-x_0},\quad z_0=0,\ z_{K-1}=1.
\label{eq:table-decomposition}
\end{equation}
Here $a=x_0$ and $R=x_{K-1}-x_0$ specify the sampled log-frequency interval,
called the support; $K-2$ interior coordinates define allocation. Every
geometric table has $z_k=k/(K-1)$ irrespective of base. The space
$\mathcal Z_K=\{0=z_0<\cdots<z_{K-1}=1\}$ includes analytic and learned allocations.

Fixing the endpoints isolates the effect of interior placement.
For deployment relative to a native table, write
\begin{equation}
\omega'_k=\omega_k^N s^{-m_k},\qquad d_k=\log(\omega_k^N/\omega'_k)=m_k\log s.
\label{eq:movement-allocation}
\end{equation}
For $m_0=0$, $m_{K-1}=1$, the corresponding allocation is
$z'_k=(R_Nz_k^N+m_k\log s)/(R_N+\log s)$, where $R_N$ is the native
log-frequency span. Thus $m$ describes placement relative to the native grid;
the distance stretch in Figure~\ref{fig:intro-allocation} is $s^{m_k}$.
Changing frequencies preserves RoPE's relative-position composition and norms
\citep{su2024roformer}. A cosine/sine gain $g$ scales the rotary Q/K contribution
by $g^2$. Complete methods may jointly choose allocation, range and amplitude.

Geo denotes each protocol's geometric reference and Native a released model's
inherited table. Standard and midpoint geometric grids share $z$ but differ by
a frequency offset (Appendix~\ref{sec:table-identities}).

\paragraph{Separating the effects of allocation.}
At fixed support, total displacement $D=\sum_kd_k$ is a statistic of $z$.
Matching $D$ tests whether the distribution of that movement matters;
permuting the same frequency multiset instead tests its assignment to learned
coordinates. Appendix~\ref{sec:identification-details}
specifies these controls.
